\section{Root profiles and the two arithmetic constructions}
\label{sec:profiles}

Write $P(x)$ for the Clebsch pentagon polynomial and $A$ for its adjacency
matrix. An added vertex has fractional adjacency proportions
$p=(p_i)_{i=1}^{16}\in[0,1]^{16}$ to sixteen unit classes. Its degree and
pentagon link are
\[
 d(p)=\sum_i p_i,\qquad
 c(p)=\sum_{\mathrm{Cl}[U]\cong P_4}
       \prod_{\substack{i\text{ endpoint}\\\text{of }U}}p_i
       \prod_{\substack{j\text{ internal}\\\text{in }U}}(1-p_j).
\]

\begin{lemma}[The profile envelope]\label{lem:profile-envelope}
For every fractional profile,
\[
 c(p)\leq\min\left\{\frac{12}{5}d(p)^2,60\right\}.
\]
For every $0\leq d\leq5$, the first bound is attained by setting
$p_i=d/5$ on a Clebsch neighborhood and zero elsewhere.
\end{lemma}
\begin{proof}
Every nonadjacent pair of Clebsch vertices is the endpoint pair of exactly
six induced four-vertex paths; adjacent pairs are endpoints of none.
Dropping the internal nonneighbor factors therefore gives
\[
 c(p)\leq6\sum_{\substack{i<j\\ij\notin E(\mathrm{Cl})}}p_i p_j.
\]
The independence number of $\mathrm{Cl}$ is five. Indeed, after translating
an independent set to contain $0$, its other vertices lie in the ten pair
vertices of the standard labeling. Pairwise nonadjacency there means
pairwise intersection of two-subsets of $[5]$. Such a family has at most
four members, since a family of size four must be a star. Clebsch
neighborhoods attain size five.

For completeness, the weighted clique inequality of Motzkin and
Straus~\cite{MotzkinStraus1965} used here follows by repeatedly merging the
weights of two nonadjacent vertices in the support of a graph. With their
sum fixed, the edge quadratic form is linear in the two weights, so one
of the two mergers does not decrease it. Eventually the support is a
clique of size $r$, where its value is at most $(1-1/r)d^2/2$ by
Cauchy--Schwarz. Apply this to the complement of $\mathrm{Cl}$, whose
clique number is five, to obtain the first bound.

The bound $c(p)\leq60$ is a finite exact ingredient: for every one of the
$65,536$ Boolean neighborhoods $N$, count the $480$ induced $P_4$ subsets
with endpoints in $N$ and internal vertices outside $N$. The supplied
standard-library program \texttt{verify\_profile\_envelope.py} constructs
these paths from the Cayley edge list and checks every count. It finds
$c(N)\leq60$, with equality precisely at the sixteen open and sixteen
closed neighborhoods. The degreewise maxima, for degrees zero through
sixteen, are
\[
 0,0,6,18,36,60,60,54,48,44,42,35,25,18,12,0,0.
\]
Multilinearity represents $c(p)$ as the expectation of $c(N)$ for an
independent Bernoulli subset, so the Boolean inequality gives the
fractional one.

Finally, a neighborhood is an independent five-set. An induced
four-vertex path with endpoints in that neighborhood has its two internal
vertices outside it. There are six paths per endpoint pair and ten pairs,
yielding $c(p)=60(d/5)^2$ for the stated profile.
\end{proof}

Define
\[
 x_m=\sqrt{\frac{32m}{5}},\qquad
 \alpha_m=x_m-\lfloor x_m\rfloor,\qquad
 k_m=\sqrt{\frac{m}{40}}=\frac{x_m}{16},
\]
and put
\[
 f(\alpha)=\frac{(1-\alpha)\min\{3\alpha,1\}}{80},
 \qquad 0\leq\alpha<1.
\]
As usual $\Delta(m)=c_\ast m^{5/2}-\rho(C_5,m)$.

\begin{proposition}[Uniform constructive upper error]
\label{prop:constructive-error}
As $m\to\infty$ through the integers,
\[
 \Delta(m)\leq m^2 f(\alpha_m)+O(m^{7/4}),
\]
with an absolute implicit constant.
\end{proposition}
\begin{proof}
Write $k=k_m$ and $\alpha=\alpha_m$. All estimates below are uniform for
$0\leq\alpha<1$.

For the first construction let $N=\lfloor x_m\rfloor$ and distribute $N$
vertices into sixteen classes whose sizes differ by at most one. Their
Clebsch blow-up $B$ has
\[
 e(B)=\frac{5N^2}{32}+O(1),\qquad
 I_5(B)=192\left(\frac{N}{16}\right)^5+O(k^3).
\]
Fix a template vertex $v$. Add a single new vertex, with neighbors only
in the five classes indexed by $N_{\mathrm{Cl}}(v)$, and choose its degree
as close to $m-e(B)$ as feasibility permits. This leaves at most $O(1)$
edges to add or delete: the available number of neighbors is $5k+O(1)$,
whereas $m-e(B)=5k\alpha+O(1)$. Within those five classes choose the
neighbor counts to be $\alpha k+O(1)$, with their total equal to the
chosen degree. Such simultaneous rounding is possible by distributing
the bounded integer remainder among the five classes.

There are sixty induced path templates with their endpoints in the
selected neighborhood. Consequently the added vertex lies in
$60\alpha^2k^4+O(k^3)$ induced pentagons. Correct the bounded edge
remainder by adding a disjoint matching or deleting edges. Each deletion
destroys at most $O(k^3)$ of the existing pentagons; adding the matching
changes none. The resulting graph has exactly $m$ edges and
\begin{align*}
 I_5(G)&=192(k-\alpha/16)^5+60\alpha^2k^4+O(k^3)\\
       &=192k^5-60\alpha(1-\alpha)k^4+O(k^3).
\end{align*}
Thus $\Delta(m)\leq60\alpha(1-\alpha)k^4+O(k^3)$.

For the second construction put $N=\lfloor x_m\rfloor+1$,
$\beta=1-\alpha$, and $t=N/16=k+\beta/16$. The Clebsch matrix has a
zero-sum vector $w\in\{-1,1\}^{16}$ satisfying $Aw=-3w$. Explicitly, at
the vertex $u\in\mathbb F_2^4$ take $w_u=(-1)^{u_1+u_2+u_3}$; it has
squared norm sixteen. Set
\[
 r^2=\frac{40t^2-m}{24},\qquad y_i=t+rw_i.
\]
Then $\sum_i y_i=N$, all $y_i$ are positive for large $m$, and
\[
 E(y)=40t^2-24r^2=m,\qquad
 \sum_i(y_i-t)^2=16r^2=\frac23(40t^2-m).
\]
Round the $y_i$ to integers $n_i$ preserving their sum $N$ and with
$|n_i-y_i|\leq1$. Since $r=O(\sqrt{k})$ and the rounding error has sum
zero, the resulting blow-up $B$ satisfies
\[
 e(B)=m+O(\sqrt{k}),\qquad
 V:=\sum_i(n_i-t)^2=\frac23(40t^2-m)+O(\sqrt{k}).
\]
The logarithmic Hessian calculation for the template, or direct
polynomial expansion, gives uniformly for zero-sum $h=O(\sqrt{k})$,
\[
 c_\ast E(t\mathbf1+h)^{5/2}-P(t\mathbf1+h)
   =6t^3\lVert h\rVert_2^2+O(k^{7/2}).
\]
Hence the deficit of $B$ at its own edge count is
\[
 6k^3\frac23(40t^2-m)+O(k^{7/2})
   =20\beta k^4+O(k^{7/2}).
\]
Correct $e(B)-m$ by adding a disjoint matching or deleting edges. The
number of changes is $O(\sqrt{k})$, and each deletion destroys at most
$O(k^3)$ pentagons. The change in the comparison expression
$c_\ast e(B)^{5/2}$ is also $O(k^{7/2})$. Thus this construction gives
$\Delta(m)\leq20(1-\alpha)k^4+O(k^{7/2})$.

Take the better construction and use $k^4=m^2/1600$.
\end{proof}

\section{A sharp arithmetic phase law}\label{sec:phase-law}

We retain $k=\sqrt{m/40}$, $x=16k=\sqrt{32m/5}$,
$\alpha=x-\lfloor x\rfloor$, and the function
\[
 f(\alpha)=\frac{(1-\alpha)\min\{3\alpha,1\}}{80}
\]
from the preceding section.

\begin{lemma}[The scalar profile problem]\label{lem:scalar-profile}
Put $\phi(u)=\max\{12u,-4u\}$. For arbitrary fractional profiles
$p_1,\ldots,p_b\in[0,1]^{16}$ and any integer $\ell$,
\begin{equation}\tag{58}\label{eq:scalar-profile}
 \sum_{j=1}^b\bigl(12d(p_j)-c(p_j)\bigr)
 +\phi\left(5(\alpha+\ell)-\sum_{j=1}^bd(p_j)\right)
 \geq20(1-\alpha)\min\{3\alpha,1\}.
\end{equation}
\end{lemma}
\begin{proof}
Write $z_j=d(p_j)/5$ and $\Phi(u)=\phi(5u)$. Lemma~11.1 gives
\[
 12d(p_j)-c(p_j)\geq g(z_j),\qquad
 g(z)=\begin{cases}
       60z(1-z),&0\leq z\leq1,\\
       60(z-1),&z\geq1.
      \end{cases}
\]
The function $\Phi$ is Lipschitz with constant sixty. Thus replacing any
$z_j>1$ by one can only decrease
$\sum_jg(z_j)+\Phi(\alpha+\ell-\sum_jz_j)$: its profile cost drops by
$60(z_j-1)$ and the other term rises by at most that amount. We may
consequently suppose $0\leq z_j\leq1$.

For a fixed sum $\sum_jz_j=a+t$, where $a$ is an integer and $0\leq t<1$,
the sum $\sum_jz_j^2$ is maximized by $a$ ones, one entry $t$, and zeros.
One can see this by moving mass between any two entries strictly between
zero and one until one reaches an endpoint. Hence
\[
 \sum_j60z_j(1-z_j)\geq60t(1-t).
\]
Absorbing $a$ into $\ell$ reduces the problem to minimizing
\[
 60t(1-t)+\Phi(\alpha+\ell-t),\qquad 0\leq t\leq1,\quad\ell\in\mathbb Z.
\]
For $\ell\leq-1$, the second term alone is at least $20(1-\alpha)$; for
$\ell\geq1$, the minimum over $t$ is at least $60\alpha$. For $\ell=0$
and $t\leq\alpha$ the expression is $60(\alpha-t^2)$, whose minimum is
$60\alpha(1-\alpha)$. For $t\geq\alpha$ it is the concave quadratic
$80t-60t^2-20\alpha$, so its minimum is attained at $t=\alpha$ or $t=1$.
The two values are $60\alpha(1-\alpha)$ and $20(1-\alpha)$, respectively.
Their minimum is the asserted right side.
\end{proof}

\begin{theorem}[Sharp phase-dependent second-order asymptotic]
\label{thm:phase-law}
As $m\to\infty$ through the positive integers,
\begin{equation}\tag{59}\label{eq:phase-law}
 \begin{gathered}
 \rho(C_5,m)=c_\ast m^{5/2}
   -\frac{m^2}{80}(1-\alpha_m)\min\{3\alpha_m,1\}+O(m^{15/8}),\\
 \alpha_m=\left\{\sqrt{\frac{32m}{5}}\right\}.
 \end{gathered}
\end{equation}
The implicit constant is absolute; in particular the remainder is uniform
as the phase approaches either endpoint.
\end{theorem}
\begin{proof}
The constructive upper bound for the deficit is Proposition~11.2. We
prove the matching lower bound with error $O(k^{15/4})=O(m^{15/8})$.
Choose an extremal graph $G$ with $m$ edges. The coarse lower construction
gives $\Delta(m)=O(m^2)=O(k^4)$, so the linear structural stability
theorem supplies a Clebsch blow-up $H$, with isolates, on the same vertex
set such that
\begin{equation}\tag{60}\label{eq:phase-structure}
 T:=|E(G)\mathbin\triangle E(H)|=O(k),\qquad
 N_H=16k+O(1),\qquad n_i^H=k+O(\sqrt{k}),
\end{equation}
and $\sum_i(n_i^H-N_H/16)^2=O(k)$. All constants in this proof are
independent of $m$ and $\alpha$.

Discard vertices isolated in both $G$ and $H$. Every remaining isolated
vertex of $H$ is incident with an edge in $E(G)\setminus E(H)$, so there
are at most $2T$ such vertices. The total number of remaining vertices is
therefore $O(k)$. This observation allows the counting estimates below to
be independent of the original number of isolates.

Let the wrong degree of a vertex be its number of incident pairs in
$E(G)\mathbin\triangle E(H)$. Call a vertex bad if its wrong degree
exceeds $k^{3/4}$. The bad set $B$ has size
\[
 b\leq2T/k^{3/4}=O(k^{1/4}).
\]
Let $G_0,H_0$ be the graphs remaining after deleting $B$. Write $n_i$ for
the sixteen active class sizes of $H_0$, put $N=\sum_i n_i$, $q=N/16$, and
$h_i=n_i-q$. Then
\begin{equation}\tag{61}\label{eq:phase-masses}
 q=k+O(k^{1/4}),\qquad n_i=k+O(\sqrt{k}),\qquad
 V:=\sum_i h_i^2=O(k).
\end{equation}
In particular every class is nonempty for all sufficiently large $m$. If
$A$ is the Clebsch adjacency matrix, put
\[
 a=\frac{h^TAh}{2k},\qquad v=V/k.
\]
The identity $A^2=3I-2A+2J$ shows that the eigenvalues on the sum-zero
subspace are $1$ and $-3$. Consequently
\begin{equation}\tag{62}\label{eq:phase-spectrum}
 -\frac32v\leq a\leq\frac12v.
\end{equation}

For a bad vertex $j$, let $p_{ji}$ be the proportion of the $i$th active
class of $H_0$ adjacent to $j$ in $G$. Denote this profile by $p_j$, and
abbreviate $d_j=d(p_j)$, $c_j=c(p_j)$. Let $t_-$ count the edges of $H_0$
missing from $G_0$, and let $t_+$ count the edges added to $H_0$ inside
$G_0$, together with the edges from $B$ to the isolated class of $H_0$.
Both are $O(k)$: the latter additional edges are all wrong pairs in the
original comparison with $H$. Let $t_n$ count those added edges inside
$G_0$ whose endpoints have distinct nonadjacent active types; thus
$0\leq t_n\leq t_+$.

The edge count is
\begin{equation}\tag{63}\label{eq:phase-edge-count}
 \frac{m-5N^2/32}{k}
 =a+\sum_{j\in B}d_j+\frac{t_+-t_-}{k}+O(k^{-1/4}).
\end{equation}
Indeed, $e(H_0)=5N^2/32+ka$. The exact number of edges from a bad
vertex to the active classes is $\sum_i n_i p_{ji}$; replacing this by
$kd_j$ has error $O(\sqrt{k})$ per bad vertex and thus $O(k^{3/4})$ in
total. There are at most $b^2=O(k^{1/2})$ edges within $B$. These are all
the remaining terms.

We next expand the pentagon count and give the needed error estimates
explicitly. Every good vertex has wrong degree at most $k^{3/4}$. If
$r_u$ denotes its wrong degree in the comparison of $G_0$ and $H_0$, then
\[
 \sum_u r_u^2\leq k^{3/4}\sum_u r_u=O(k^{7/4}).
\]
An intersecting pair of wrong edges fixes three vertices and belongs to
$O(k^2)$ five-sets; a disjoint pair fixes four vertices and belongs to
$O(k)$ five-sets. The number of five-sets containing at least two wrong
edges is therefore
\begin{equation}\tag{64}\label{eq:wrong-pair-error}
 O\left(k^2\sum_u r_u^2+kT^2\right)=O(k^{15/4}).
\end{equation}
For a five-set with at most one wrong pair, the change of its
induced-pentagon indicator is exactly its one-edge change. For every
other five-set, the difference between the actual change and the sum of
its one-edge changes has absolute value at most eleven. Hence (64)
bounds the entire higher-order error.

By the two-vertex coefficients of Lemma~A.3, one deleted template edge
loses $24k^3+O(k^{5/2})$ pentagons, and an added edge between distinct
nonadjacent active types has change $-6k^3+O(k^{5/2})$. Edges added
within a class or meeting an isolated type have one-edge change zero.
The estimates are uniform because of (61); excluding the two fixed
endpoints from the available old classes introduces only $O(k^2)$
further error. There are $O(k)$ wrong edges. Thus
\begin{equation}\tag{65}\label{eq:good-pentagons}
 I_5(G_0)=P(n)-24k^3t_--6k^3t_n+O(k^{15/4}).
\end{equation}

Pentagons containing at least two bad vertices number at most
$O(b^2k^3)=O(k^{7/2})$. For a fixed bad vertex, changing the background
on the other four vertices from $G_0$ to $H_0$ affects at most
$O(Tk^2)$ five-sets. Summing over $B$ costs $O(bTk^2)=O(k^{13/4})$.
In the perfect background $H_0$, the root link is obtained from $c(p_j)$
by giving its four distinct classes weights $n_i$. Each product of four
such weights is $k^4+O(k^{7/2})$. There are only $480$ path templates, so
this link is $k^4c_j+O(k^{7/2})$, uniformly in the profile. Summing the
latter errors over $B$ gives $O(k^{15/4})$. Isolated background vertices
lie in no such path. We conclude from (65) that
\begin{equation}\tag{66}\label{eq:phase-pentagon-count}
 I_5(G)=P(n)+k^4\sum_{j\in B}c_j-24k^3t_--6k^3t_n+O(k^{15/4}).
\end{equation}

It remains to expand the smooth template defect. For
$Q(x)=c_\ast E(x)^{5/2}-P(x)$, the Clebsch counts yield
\[
 Q(\mathbf1)=0,\qquad\nabla Q(\mathbf1)=0,\qquad
 \nabla^2Q(\mathbf1)=12I-\tfrac34J.
\]
Homogeneity, $\sum_i h_i=0$, and Taylor's theorem give
\[
 Q(n)=6q^3V+O(q^2\lVert h\rVert_2^3)=6k^3V+O(k^{7/2}).
\]
Also $e(H_0)=m+O(k^{5/4})$, by (61) and $V=O(k)$. Since
$(5/2)c_\ast m^{3/2}=12k^3$, another Taylor expansion gives
\[
 c_\ast m^{5/2}-c_\ast e(H_0)^{5/2}
   =12k^3\bigl(m-e(H_0)\bigr)+O(k^{7/2}).
\]
Combining these identities with (63) and (66) produces
\begin{equation}\tag{67}\label{eq:phase-deficit}
 \frac{\Delta(m)}{k^4}
 =6v+\sum_{j\in B}(12d_j-c_j)
    +12\frac{t_++t_-}{k}+6\frac{t_n}{k}+O(k^{-1/4}).
\end{equation}

The scalar function $\phi$ is subadditive and Lipschitz with constant
twelve. Equation (62) implies $6v\geq\phi(a)$, while
\[
 12(t_++t_-)/k\geq\phi\bigl((t_+-t_-)/k\bigr).
\]
Drop the nonnegative $t_n$ term and apply (63). If
$r=(m-5N^2/32)/k$, we obtain
\[
 \frac{\Delta(m)}{k^4}
 \geq\sum_{j\in B}(12d_j-c_j)
    +\phi\left(r-\sum_{j\in B}d_j\right)-O(k^{-1/4}).
\]
Since $N=16k+O(k^{1/4})$,
\begin{align*}
 r&=5(16k-N)-\frac{5}{32k}(16k-N)^2\\
  &=5(\alpha+\ell)+O(k^{-1/2}),\qquad
       \ell=\lfloor16k\rfloor-N\in\mathbb Z.
\end{align*}
Apply Lemma~12.1 and use the Lipschitz bound once more. The result is
\[
 \Delta(m)\geq20k^4(1-\alpha)\min\{3\alpha,1\}-O(k^{15/4}).
\]
Because $m^2=1600k^4$, this is the desired lower bound.
Proposition~11.2 has an upper error $O(m^{7/4})$, which is absorbed in
$O(m^{15/8})$.
\end{proof}

\begin{corollary}[The optimal uniform second-order coefficient]
\label{cor:uniform-coefficient}
The finite-edge deficit satisfies
\[
 \limsup_{m\to\infty}\frac{\Delta(m)}{m^2}=\frac1{120},\qquad
 \liminf_{m\to\infty}\frac{\Delta(m)}{m^2}=0.
\]
Its set of subsequential limits is the whole interval $[0,1/120]$. In
particular, the coefficient $1/120$ in the uniform estimate
\[
 \rho(C_5,m)\geq c_\ast m^{5/2}-\frac{m^2}{120}-O(m^{15/8})
\]
cannot be decreased.
\end{corollary}
\begin{proof}
The function $f$ is continuous on the unit circle, vanishes at phase zero,
and has maximum $1/120$ at phase $1/3$. For each fixed $0<\beta<1$,
choose an integer $m_k$ nearest to $(5/32)(16k+\beta)^2$. Factoring the
difference of squares gives
$\sqrt{32m_k/5}=16k+\beta+O(k^{-1})$, so $\alpha_{m_k}\to\beta$.
The exact equality sequence $m=40k^2$ provides phase zero. The theorem
therefore yields all values of the continuous range
$f([0,1])=[0,1/120]$ and excludes all other limit points.
\end{proof}

\section{The lattice correction at a fixed edge offset}
\label{sec:lattice-offset}

The phase law describes the $m^2$ scale. In a bounded edge neighborhood
of the equality sequence, the deficit instead has order $m^{3/2}$. Its
leading coefficient is a finite lattice optimization determined by the
Clebsch adjacency matrix.

For $h\in\mathbb Z^{16}$ with $\sum_i h_i=0$, put
\[
 Q(h)=\sum_{uv\in E(\mathrm{Cl})}h_uh_v=\tfrac12h^TAh.
\]
For an integer $j$, define
\begin{equation}\tag{68}\label{eq:lattice-objective}
 \Lambda(j)=\min_{\substack{h\in\mathbb Z^{16}\\\sum_i h_i=0}}
 \bigl(6\lVert h\rVert_2^2+12|j-Q(h)|\bigr).
\end{equation}
The minimum may be restricted to $\lVert h\rVert_2^2\leq2|j|$: the zero
vector gives value $12|j|$, and every vector outside that ball has
strictly greater cost. In particular this is a finite, explicitly
decidable optimization for each $j$.

\begin{theorem}[Fixed-offset asymptotic]\label{thm:fixed-offset}
For every fixed integer $j$, as the positive integer $k$ tends to infinity,
\begin{equation}\tag{69}\label{eq:fixed-offset-asymptotic}
 \rho(C_5,40k^2+j)
   =c_\ast(40k^2+j)^{5/2}-\Lambda(j)k^3+O_j(k^2).
\end{equation}
The coefficients satisfy
\[
 \Lambda(j)=12j\quad(j\geq0),\qquad
 4|j|\leq\Lambda(j)\leq12|j|\quad(j<0),
\]
and $\Lambda(-12)=48$. Thus the fixed-offset correction is not symmetric
under $j\mapsto-j$.
\end{theorem}
\begin{proof}
Write $m=40k^2+j$ and $\Delta(m)=c_\ast m^{5/2}-\rho(C_5,m)$. Start
from the balanced blow-up of class size $k$. If $j\geq0$, add a disjoint
matching with $j$ edges; if $j<0$, delete $|j|$ edges. Since a single
deletion affects $O(k^3)$ pentagons, this gives the preliminary estimate
$\Delta(m)=O_j(k^3)$.

Let $G$ be extremal with $m$ edges. Its normalized deficit is
$O_j(k^{-2})$, so Theorem~9.3 supplies a Clebsch blow-up $H$, with
isolates, for which
\[
 |E(G)\mathbin\triangle E(H)|=O_j(1),\qquad
 \sum_i(n_i-N/16)^2=O_j(1),
\]
\[
 N=\sqrt{32m/5}+O_j(k^{-1})=16k+O_j(k^{-1}).
\]
The integer $N$ therefore equals $16k$ for all sufficiently large $k$.
Consequently $n_i=k+h_i$ for an integer vector $h$ satisfying
$\sum_i h_i=0$ and $\lVert h\rVert_2=O_j(1)$, and
\begin{equation}\tag{70}\label{eq:offset-edge-count}
 e(H)=40k^2+Q(h).
\end{equation}
After discarding common isolates, the two graphs have $O(k)$ vertices:
every remaining isolated vertex of $H$ is incident with one of the
$O_j(1)$ wrong edges.

Let $t_+$ and $t_-$ be the numbers of added and deleted edges in passing
from $H$ to $G$, and let $t_n$ count added edges between distinct
nonadjacent active classes. All these numbers are $O_j(1)$, and
\[
 j-Q(h)=t_+-t_-.
\]
The one-edge pentagon changes are $-24k^3+O_j(k^2)$ for a deleted
template edge, $-6k^3+O_j(k^2)$ for an added edge of the latter
nonadjacent type, and zero for an edge within a class or meeting an
isolated type. Five-sets containing at least two wrong pairs number
$O_j(k^2)$, by the same pair counting used in (64). Therefore
\[
 I_5(G)=P(k\mathbf1+h)-24k^3t_--6k^3t_n+O_j(k^2).
\]
The template Hessian gives
\[
 c_\ast e(H)^{5/2}-P(k\mathbf1+h)
   =6k^3\lVert h\rVert_2^2+O_j(k^2).
\]
Since $m-e(H)=j-Q(h)=O_j(1)$, the comparison between the two edge
powers is
\[
 c_\ast m^{5/2}-c_\ast e(H)^{5/2}
   =12k^3\bigl(j-Q(h)\bigr)+O_j(k).
\]
Combining these formulas yields
\begin{align*}
 \Delta(m)
 &=k^3\bigl(6\lVert h\rVert_2^2+12(t_++t_-)+6t_n\bigr)+O_j(k^2)\\
 &\geq k^3\bigl(6\lVert h\rVert_2^2+12|j-Q(h)|\bigr)-O_j(k^2)\\
 &\geq\Lambda(j)k^3-O_j(k^2).
\end{align*}

For the reverse inequality, fix a minimizing vector $h$ in (68). The
class sizes $k+h_i$ are positive for large $k$. If $j-Q(h)>0$, add that
many edges within a class; if it is negative, delete its absolute value
in template edges. There is ample capacity for either operation once
$k$ is large. All added edges are within a class, so $t_n=0$, and the
number of edits is exactly $|j-Q(h)|$. The same expansion gives a graph
with exactly $m$ edges whose deficit is $\Lambda(j)k^3+O_j(k^2)$. This
proves (69).

On the sum-zero subspace, the eigenvalues of $A$ are $1$ and $-3$. Hence
\[
 -\tfrac32\lVert h\rVert_2^2\leq Q(h)\leq\tfrac12\lVert h\rVert_2^2.
\]
For $j\geq0$ this implies
\[
 6\lVert h\rVert_2^2+12|j-Q(h)|
 \geq12Q(h)+12|j-Q(h)|\geq12j,
\]
and $h=0$ attains equality. For $j<0$, it implies
\[
 6\lVert h\rVert_2^2+12|j-Q(h)|
 \geq-4Q(h)+4|j-Q(h)|\geq4|j|.
\]
The upper bound $12|j|$ again follows from $h=0$.

To evaluate the example $j=-12$, use the Cayley labeling
$u\in\mathbb F_2^4$ and the characters $\chi_a(u)=(-1)^{a\cdot u}$.
For the bit masks $a=7$ and $a=11$, each character has Clebsch eigenvalue
$-3$. The vector
\[
 h=\tfrac12(\chi_7+\chi_{11})
\]
is integral, has sum zero, has eight nonzero entries in $\{-1,1\}$, and
satisfies $Ah=-3h$. Thus $\lVert h\rVert_2^2=8$, $Q(h)=-12$, and its
objective value is $48$. The lower bound $4|j|$ proves that this value is
optimal.
\end{proof}

\begin{corollary}[Two exact finite-edge values]
\label{cor:two-exact-offsets}
For all sufficiently large integers $k$,
\[
 \rho(C_5,40k^2+1)=192k^5,\qquad
 \rho(C_5,40k^2-1)=192k^5-24k^3.
\]
Up to isolated vertices, all extremal graphs at $40k^2+1$ edges are
obtained from the balanced Clebsch blow-up of class size $k$ by adding
one edge either within a class or incident with an isolated vertex. At
$40k^2-1$ edges, all extremal graphs are either a balanced blow-up with
one edge deleted, or a Clebsch blow-up in which two adjacent classes
have sizes $k+1,k-1$ and the other fourteen classes have size $k$.
\end{corollary}
\begin{proof}
For $j=1$, the lattice minimum is $12$ by the theorem. For $j=-1$, it is
also $12$: the zero vector attains this value, and every nonzero integral
sum-zero vector has squared norm at least two, so its first cost term
alone is at least twelve.

For an extremal graph and its template from the proof above, write
\[
 L=6V+12(t_++t_-)+6t_n.
\]
This is an integer multiple of six and is at least twelve. The same
proof gives $\Delta(m)=Lk^3+O(k^2)$, with a constant uniform over the
bounded family of permitted $h$ and wrong-edge counts. The elementary
constructions for $j=\pm1$ have deficit $12k^3+O(k)$. Hence $L>12$,
which would imply $L\geq18$, is impossible for all sufficiently large
$k$.

The nonnegative integer solutions to $L=12$ are therefore
$V=0$, $t_++t_-=1$, $t_n=0$, or $V=2$, $t_++t_-=t_n=0$. Here we used
that an integer sum-zero vector has even squared norm. In the second
case $h$ has one entry $1$, one entry $-1$, and all others zero, so
$Q(h)$ is zero or $-1$.

For $j=1$, only the first case is possible. The single modification is
an added edge, and $t_n=0$ says that it lies within a class or meets an
isolated vertex. Such an edge changes no pentagon count: in the first
case its endpoints are true twins, and in the other cases any formerly
isolated endpoint has degree one. This gives the asserted value and
classification.

For $j=-1$, the first case deletes one edge of the balanced template.
Exactly $24k^3$ old pentagons use that edge. There are no newly induced
pentagons: the deleted endpoints have no common neighbor, whereas a
nonadjacent pair in a five-cycle must have one. In the second case
$Q(h)=-1$ forces the two perturbed classes to be adjacent.
Multilinearity of the pentagon polynomial gives exactly
\[
 P(k\mathbf1+h)=192k^5+60k^4\sum_i h_i-24k^3=192k^5-24k^3.
\]
These constructions attain the common value, completing both the upper
bound and the equality classification.
\end{proof}

\section{Eventual polynomiality at every fixed offset}
\label{sec:eventual-polynomial}

The lattice coefficient has a further consequence. For each fixed edge
offset, the exact extremal function is eventually a polynomial, and the
polynomial can be determined by a finite enumeration. This statement
does not provide an explicit threshold beyond which the polynomial is
valid.

\begin{theorem}[Computable eventual polynomial]
\label{thm:eventual-polynomial}
For every integer $j$, there are computable integers $a_j,b_j,c_j$ and
an integer $k_0(j)$ such that, for every integer $k\geq k_0(j)$,
\begin{equation}\tag{71}\label{eq:eventual-polynomial}
 \rho(C_5,40k^2+j)
 =192k^5+\bigl(12j-\Lambda(j)\bigr)k^3+a_jk^2+b_jk+c_j.
\end{equation}
Moreover, a finite procedure produces a finite list of graph
constructions that, for all sufficiently large $k$, contains all
extremal graphs up to isomorphism and the addition of isolates. The
proof determines the eventual polynomial and this list; it does not
give an explicit value or an algorithm for $k_0(j)$.
\end{theorem}
\begin{proof}
We first strengthen the discrete conclusion in the fixed-offset
argument. Its template has integer class sizes $k+h_i$ with
$\sum_i h_i=0$ and bounded $V=\lVert h\rVert_2^2$, and a bounded number
$T=t_++t_-$ of wrong pairs. As in that proof, put
\[
 L=6V+12T+6t_n.
\]
The edge-offset constraint and the definition of $\Lambda(j)$ give
$L\geq\Lambda(j)$. Both numbers are integer multiples of six. The same
uniformly bounded-template expansion gives
\[
 \Delta(40k^2+j)=Lk^3+O_j(k^2).
\]
On the other hand, the construction in Theorem~13.1 gives
$\Delta(40k^2+j)\leq\Lambda(j)k^3+O_j(k^2)$. The positive gap between
distinct multiples of six therefore forces
\begin{equation}\tag{72}\label{eq:energy-equality}
 L=\Lambda(j)
\end{equation}
for every extremal graph and all sufficiently large $k$. Since
$\Lambda(j)\leq12|j|$, this implies the explicit bounds
\begin{equation}\tag{73}\label{eq:explicit-discrete-bounds}
 V\leq2|j|,\qquad T\leq|j|.
\end{equation}
The initial boundedness and the uniform error in this argument come
from strong stability, before imposing the explicit bounds (73); thus
there is no circular assumption about the extremal family.

We now describe a finite list depending only on $j$. Enumerate:
\begin{enumerate}
\item integer vectors $h\in\mathbb Z^{16}$ with $\sum_i h_i=0$ and
      $\lVert h\rVert_2^2\leq2|j|$;
\item a set $X$ of $s\leq2|j|$ labeled exceptional vertices, each
      assigned a type $\tau(x)\in\{0,1,\ldots,16\}$;
\item a set $F\subseteq\binom{X}{2}$ of at most $|j|$ wrong pairs, with
      every vertex of $X$ incident with a member of $F$.
\end{enumerate}
Type zero is the isolated type. The case $F=X=\varnothing$ is allowed.
The template adjacency between two exceptional vertices is Clebsch
adjacency between their positive types, and is zero if either type is
zero. A wrong pair toggles this template adjacency. Thus $F$ determines
$t_+,t_-$, and $t_n$ without any further choices. Retain precisely the
data satisfying
\begin{equation}\tag{74}\label{eq:admissible-data}
 Q(h)+t_+-t_-=j,\qquad
 6\lVert h\rVert_2^2+12|F|+6t_n=\Lambda(j).
\end{equation}
All these tests use integer arithmetic, and the enumeration is finite
with explicit bounds.

Each retained datum defines a graph for all sufficiently large $k$. Let
$e_i$ be the number of exceptional vertices of type $i$, and put
$k+h_i-e_i$ ordinary vertices in each active class $i$. Together with
the exceptional vertices, the active class sizes of the underlying
template are exactly $k+h_i$. Include the $e_0$ exceptional vertices of
type zero, place all template edges, and toggle the pairs in $F$. The
counts $k+h_i-e_i$ are nonnegative for sufficiently large $k$; this
capacity requirement has an explicit threshold for each datum. The
edge count is exactly $40k^2+j$ by (74). Further isolates may be added
freely.

Conversely, every extremal graph for sufficiently large $k$ has such a
representation. Take the endpoints of all its wrong pairs relative to
the strong-stability template as $X$. There are at most $2T\leq2|j|$ of
them. Outside $X$, all adjacencies agree with the template, and the
remaining type-zero vertices are isolated. The bounds and both
admissibility tests follow from (73), the exact edge count, and (72).

For each retained datum $D$, its pentagon count is an integer
coefficient polynomial $P_D(k)$. To see this without an automorphism
division, take a subset $S\subseteq X$ of exceptional vertices and a
subset $I\subseteq[16]$ of ordinary types, with $|S|+|I|=5$. Use one
formal ordinary representative of every type in $I$, and test whether
the resulting five-vertex graph is an induced pentagon. Let its
indicator be $b_D(S,I)$. All ordinary vertices of the same type are
false twins, even in their adjacencies to $X$, since wrong pairs have
both endpoints in $X$. An induced pentagon cannot contain two such
vertices. It follows that the exact count is
\begin{equation}\tag{75}\label{eq:exact-construction-polynomial}
 P_D(k)=\sum_{\substack{S\subseteq X,\ I\subseteq[16]\\|S|+|I|=5}}
              b_D(S,I)\prod_{i\in I}(k+h_i-e_i).
\end{equation}
Each five-set is counted once. This formula establishes both
integrality of the coefficients and an explicit procedure to compute
them.

Every retained datum has the same three leading coefficients. Indeed,
the fixed-offset expansion applies to each of these fixed finite
modifications and gives
\[
 c_\ast(40k^2+j)^{5/2}-P_D(k)=\Lambda(j)k^3+O_j(k^2).
\]
Since
\[
 c_\ast(40k^2+j)^{5/2}=192k^5+12jk^3+O_j(k),
\]
the polynomial $P_D$ has coefficient $192$ at degree five, zero at
degree four, and $12j-\Lambda(j)$ at degree three. Write its remaining
coefficients as $(a_D,b_D,c_D)$.

The retained family is nonempty. A minimizing lattice vector in (68),
followed by the addition of within-class edges or deletion of template
edges as in the fixed-offset proof, supplies such a datum. The wrong
pairs may be chosen disjoint; their number is at most $|j|$, and their
endpoints satisfy the stated bound.

For all sufficiently large $k$, every retained graph is feasible, and
every extremal graph is represented. Thus
\[
 \rho(C_5,40k^2+j)=\max_D P_D(k).
\]
Among finitely many polynomials, the one with lexicographically largest
coefficient vector, read from highest degree to lowest, is maximal for
all sufficiently large arguments. Choose the largest triple
$(a_D,b_D,c_D)$ and denote it by $(a_j,b_j,c_j)$. This proves (71).
Retaining all data whose polynomial equals that maximum gives exactly
the eventual extremal constructions, allowing duplicate isomorphic
presentations. For $j=0$ the list consists just of the balanced
template, giving $192k^5$ as required.

Finally, the finite polynomial comparisons and the capacity thresholds
can be computed from the list. The proof that all extremal graphs
enter that list uses the qualitative compactness step in strong
stability. We have not quantified its threshold, so the present
procedure computes the eventual polynomial and the eventual
construction list, rather than certifying their validity at an
arbitrary supplied finite value of $k$.
\end{proof}

\subsection{Explicit eventual polynomials at small offsets}

The finite procedure can be carried out directly for the first five
offsets. Its results, valid for all sufficiently large integers $k$, are
\begin{center}
\begin{tabular}{rrl}
\toprule
$j$&$\Lambda(j)$&$\rho(C_5,40k^2+j)$\\
\midrule
$-2$&$24$&$192k^5-48k^3+6k^2$\\
$-1$&$12$&$192k^5-24k^3$\\
$0$&$0$&$192k^5$\\
$1$&$12$&$192k^5$\\
$2$&$24$&$192k^5+6k^2$\\
\bottomrule
\end{tabular}
\end{center}
The cases $j=0,\pm1$ were proved above. We explain both the finite
check and matching constructions for $j=\pm2$.

For both signs, the lattice energy is $24$. For $j=2$ this follows from
Theorem~13.1. For $j=-2$, a nonzero sum-zero integer vector has $V\geq2$;
if $V=2$, its edge quadratic form is zero or $-1$, and its total lattice
cost is at least $24$. Vectors with $V\geq4$ have first cost term at
least $24$, while the zero vector attains $24$.

At an extremum the energy equality therefore leaves only
\[
 (V,T)=(0,2),(2,1),(4,0),\qquad t_n=0.
\]
For $V=2$ the class deviation vector has one entry $1$ and one entry
$-1$. For $V=4$ it has two entries of each sign. Substituting these
possibilities in the exact polynomial formula (75) gives only the
corrections
\[
 \{0,6k^2\}
\]
relative to $192k^5$ when $j=2$. When $j=-2$, the corrections relative
to $192k^5-48k^3$ are precisely
\[
 \{0,2k,6k^2,-6k^2,-2k^2+2k,2k^2+2k,-2k^2,2k^2\}.
\]
The enumeration uses the following convention. The sixteen Clebsch types
remain labeled, and the deviation vector $h$ is not identified under
Clebsch automorphisms. Exceptional vertices are identified only under
isomorphisms of their wrong-pair graph $F$ that preserve their types.
Thus a single wrong pair has unordered endpoint types, a $P_3$ has a
fixed center and unordered leaf types, and $2K_2$ has two unordered
endpoint pairs that are themselves unordered. Repeated types are allowed;
the exceptional vertices themselves are distinct. This convention gives
\begin{center}
\begin{tabular}{crr}
\toprule
$(V,T)$ & $j=2$ & $j=-2$\\
\midrule
$(0,2)$ & $762$ & $1,060$\\
$(2,1)$ & $0$ & $3,200$\\
$(4,0)$ & $120$ & $1,200$\\
\midrule
Total & $882$ & $5,460$\\
\bottomrule
\end{tabular}
\end{center}

Here is a direct derivation of every entry. An allowed added edge has
both endpoints in one active class or at least one endpoint of isolated
type. There are $16+17=33$ unordered type pairs of this kind. With $V=0$
and $j=2$, two disjoint additions therefore give $\binom{34}{2}=561$
presentations. For incident additions, an isolated center gives
$\binom{18}{2}=153$ unordered leaf-type pairs. An active center has only
its own type and the isolated type available for either leaf, giving
$16\binom{3}{2}=48$ presentations. The first total is consequently
$561+153+48=762$.
For $j=-2$, a deleted edge has one of the $40$ Clebsch edge types.
Two disjoint deletions give $\binom{41}{2}=820$ presentations, while
incident deletions give $16\binom{6}{2}=240$. Their sum is $1,060$.

For $V=2$, write $h=e_u-e_v$, so $Q(h)=-A_{uv}$ is zero or $-1$.
One modification cannot then produce offset $2$. To produce offset $-2$,
the ordered pair $(u,v)$ must be adjacent and the modification must be a
deletion. There are $80\cdot40=3,200$ choices.

For $V=4$, the support of $h$ consists of two positive and two negative
vertices. The Clebsch graph has $60$ induced copies of $2K_2$ and $480$
of $P_3+K_1$. To verify these counts directly, the six common nonneighbors
of the endpoints of an edge induce a matching of size three: each has
four neighbors in the union of the endpoint neighborhoods and hence one
among those six vertices. Thus the first count is $40\cdot3/2=60$.
There are $16\binom{5}{2}=160$ induced paths of order three. The union
of the neighborhoods of their three vertices has size $13$, so each path
has exactly three possible isolated fourth vertices. This gives
$160\cdot3=480$.
The condition $Q(h)=2$ requires two edges within the sign classes and no
cross edge, giving $2\cdot60=120$ sign assignments. The condition
$Q(h)=-2$ requires exactly two cross edges and no edge within a sign
class: a within-class edge together with the required three or more cross
edges would create a triangle. The two cross edges are either disjoint,
with four balanced sign assignments, or form a path with an isolated
vertex, with two. Hence the last negative-offset count is
$4\cdot60+2\cdot480=1,200$.

The supplied integer-arithmetic programs independently check this branch
census and evaluate (75) on every presentation. The two offsets give
respectively two and eight distinct polynomials, exactly as listed above.
For $j=-1,0,1$ the same convention gives $120,1,33$ presentations;
there are $6,496$ in total over the five offsets. In both displayed
correction lists for $j=\pm2$, the maximum is $6k^2$ for every positive
$k$. Eventual global extremality uses Theorem~14.1, not merely this finite
comparison.

For an explicit construction at $j=2$, add a new vertex to the balanced
blow-up and join it to two old vertices in distinct nonadjacent Clebsch
classes. There are six induced path templates between those types, and
the two internal vertices can each be chosen in $k$ ways. Thus exactly
$6k^2$ induced pentagons are added. At $j=-2$, delete two edges with a
common endpoint and their other endpoints in distinct classes. Each
deleted edge lies in $24k^3$ pentagons, and exactly $6k^2$ pentagons
contain both. Deletions create no new induced pentagons in the
triangle-free template. This realizes $192k^5-48k^3+6k^2$ and completes
the two examples.
